\section{Conclusiones}
\label{sec:conclusiones}
La integración final reúne 21\,600 RTT válidos en 112 corridas,
después de excluir 4\,800 observaciones de calentamiento. Los datos
se mantienen separados en tres equipos con mediciones de ambos
mecanismos. Las condiciones secuenciales tienen cinco repeticiones;
gRPC concurrente de Steven conserva una por tamaño.

La dirección de la hipótesis se observó en los tres equipos para
100 bytes y envíos secuenciales: las medianas de sockets/gRPC fueron
0,0901/1,28585 ms en Melanir, 0,0479/0,33505 ms en Mourad y
0,0231/0,27715 ms en DESKTOP-5MJCPHU. En Melanir, la mediana de
sockets fue un 92,99\,\% menor respecto de gRPC; este porcentaje
no se extrapola a las otras máquinas ni a redes físicas.

Con 10\,000 bytes y tres clientes en Mourad, las medias fueron
0,1639 ms en sockets y 0,6005 ms en gRPC. Los resultados describen
las implementaciones ensayadas: no aíslan el transporte, porque
la serialización y los intervalos de cronometraje difieren. Las
tasas concurrentes de Steven requieren especial cautela por su
distinto número de corridas y la preparación de canales excluida
del intervalo gRPC.

Las series iniciales de Lamport verificaron 21 relaciones y 42
eventos por corrida. Tras incorporar los acuses, la sesión gRPC de
Steven del 21/11 comprobó 42 relaciones causales y 84 actualizaciones
locales correctas. Dos envíos de nodos distintos con reloj 1 carecen
de camino causal en esos registros: Lamport ordena la precedencia
registrada, pero no asigna una marca única a cada evento.

Para el SGIRI documentado por Angel, gRPC aporta un contrato tipado
en la consulta de incidencias al servicio de técnicos; sockets es
una opción para un canal interno simple de estado de equipos de red.
Estas son elecciones de diseño sustentadas en el contrato y el tipo
de interacción, no pruebas de rendimiento o tolerancia a fallos del
PFC. El grupo conserva código, CSV, hashes y evidencias para reproducir
el análisis y sustentar sus aportes.
